\providecommand{\K}{\mathbb{K}}

\subsection{Формулировка теоремы Банаха об обратном операторе.
	Доказательство теоремы Банаха в случае гильбертова пространства}

\begin{theorem}[Банаха об обратном операторе]
	Пусть \(E_1,E_2\) --- банаховы пространства,
	\[
	A\in\mathcal{L}(E_1,E_2),
	\]
	и \(A:E_1\to E_2\) --- биекция. Тогда
	\[
	A^{-1}\in\mathcal{L}(E_2,E_1).
	\]
\end{theorem}

\begin{theorem}[Гильбертов случай]
	Пусть \(H\) --- гильбертово пространство,
	\[
	A\in\mathcal{L}(H),
	\]
	и \(A:H\to H\) --- биекция. Тогда
	\[
	A^{-1}\in\mathcal{L}(H).
	\]
\end{theorem}

\begin{proof}
	Сначала получим ограниченный обратный к \(A^*\), а затем применим то же рассуждение
	к оператору \(A^*\).
	
	\textbf{1. Докажем, что \(A^*\) ограничен снизу.}
	
	По билету 8.1 для этого достаточно получить оценку
	\[
	\exists m>0
	\quad
	\forall x\in H
	\qquad
	\|A^*x\|\geqslant m\|x\|.
	\]
	
	Оценка однородна по \(x\), поэтому достаточно рассмотреть векторы, для которых
	\[
	\|A^*x\|=1.
	\]
	Положим
	\[
	S=\{x\in H:\|A^*x\|=1\}.
	\]
	Тогда достаточно доказать, что \(S\) ограничено:
	\[
	\exists C>0
	\quad
	\forall x\in S
	\qquad
	\|x\|\leqslant C.
	\]
	
	По теореме Хана достаточно показать, что \(S\) слабо ограничено.
	
	Пусть \(y\in H\). Так как \(A\) сюръективен, существует \(z\in H\), такое что
	\[
	Az=y.
	\]
	Тогда для любого \(x\in S\)
	\[
	|(x,y)|
	=
	|(x,Az)|
	=
	|(A^*x,z)|
	\leqslant
	\|A^*x\|\|z\|
	=
	\|z\|.
	\]
	Значит, \(S\) слабо ограничено, а по теореме Хана --- ограничено:
	\[
	\exists C>0
	\quad
	\forall x\in S
	\qquad
	\|x\|\leqslant C.
	\]
	
	Теперь перейдём от ограниченности \(S\) к оценке снизу.
	Для \(u\neq0\) сначала проверим, что \(A^*u\neq0\).
	
	Так как \(A\) сюръективен,
	\[
	\operatorname{Im}A=H,
	\]
	поэтому
	\[
	\ker A^*
	=
	(\operatorname{Im}A)^\perp
	=
	\{0\}.
	\]
	Следовательно,
	\[
	\frac{u}{\|A^*u\|}\in S.
	\]
	Отсюда
	\[
	\frac{\|u\|}{\|A^*u\|}
	\leqslant C,
	\]
	то есть
	\[
	\|A^*u\|
	\geqslant
	\frac1C\|u\|.
	\]
	Значит, \(A^*\) ограничен снизу.
	
	\textbf{2. Покажем, что \(A^*\) --- биекция.}
	
	Инъективность уже получена:
	\[
	\ker A^*=\{0\}.
	\]
	
	Осталось доказать сюръективность. Так как \(A\) инъективен,
	\[
	\ker A=\{0\}.
	\]
	По разложению
	\[
	H=\overline{\operatorname{Im}A^*}\oplus\ker A
	\]
	получаем
	\[
	\overline{\operatorname{Im}A^*}=H.
	\]
	
	С другой стороны, \(A^*\) ограничен снизу, поэтому по билету 8.1
	\[
	\operatorname{Im}A^*
	\]
	замкнут. Следовательно,
	\[
	\operatorname{Im}A^*
	=
	\overline{\operatorname{Im}A^*}
	=
	H.
	\]
	
	Таким образом, \(A^*\) --- биекция. Так как \(A^*\) ограничен снизу,
	по билету 8.1
	\[
	(A^*)^{-1}\in\mathcal L(H).
	\]
	
	\textbf{3. Применим доказанное к \(A^*\).}
	
	Оператор
	\[
	A^*\in\mathcal L(H)
	\]
	является биекцией. Поэтому, повторяя доказанное выше для \(A^*\), получаем
	\[
	((A^*)^*)^{-1}\in\mathcal L(H).
	\]
	Так как
	\[
	(A^*)^*=A,
	\]
	имеем
	\[
	\boxed{A^{-1}\in\mathcal L(H)}.
	\]
\end{proof}
